\documentclass[11pt,a4paper]{article}

\usepackage[utf8]{inputenc}
\usepackage[T1]{fontenc}
\usepackage{lmodern}
\usepackage[margin=2.4cm]{geometry}
\usepackage{amsmath,amssymb,bm}
\usepackage{booktabs}
\usepackage{longtable}
\usepackage{array}
\usepackage{enumitem}
\usepackage{xcolor}
\usepackage{fancyvrb}
\usepackage[hidelinks]{hyperref}

\renewcommand{\tablename}{Tabel}
\renewcommand{\contentsname}{Daftar Isi}
\renewcommand{\refname}{Rujukan Model}
\setlength{\parskip}{0.45em}
\setlength{\parindent}{0pt}

\newcommand{\nhat}{\hat{\bm{n}}}
\newcommand{\shat}{\hat{\bm{s}}}
\newcommand{\what}{\hat{\bm{w}}}
\newcommand{\xp}{\hat{\bm{x}}_p}
\newcommand{\yp}{\hat{\bm{y}}_p}
\newcommand{\degg}{^\circ}
\newcommand{\model}[1]{\textbf{\textit{Model: #1}}}

\title{\textbf{Solar Tracker Dua Sumbu PENS\\ sebagai Cyber-Physical System}\\[0.4em]
\large Model Matematis, Validasi Numerik, dan Jalan Menuju Digital Twin}
\author{Program Magister Terapan, Politeknik Elektronika Negeri Surabaya\\
Titik pemasangan: $-7{,}276006\degg$ LS, $112{,}794327\degg$ BT}
\date{\today}

\begin{document}
\maketitle

\begin{abstract}
\noindent
Dokumen ini memodelkan solar tracker dua sumbu (azimuth dan kemiringan) sebagai Cyber-Physical System,
mengikuti kerangka \emph{Solar Tracking Systems and Cyber-Physical Systems}, tetapi dengan seluruh besaran
berasal dari CAD alat, pelat nama komponen, data cuaca terukur, atau simulasi multi-benda yang tervalidasi.
Setiap persamaan diberi nama modelnya dan setiap asumsi diberi tanda eksplisit.

Alat yang dimodelkan memakai \textbf{motor stepper loop tertutup 57HSE1.5N} ($T_h = 1{,}5$~N\,m) dengan
enkoder di poros rotor, dan \textbf{reduksi gigi $5{,}0{:}1$ pada sumbu kemiringan} (pinion 6 gigi, gigi
besar 30, modul $2{,}75$; rasio diukur langsung di alat). Sumbu azimuth tetap penggerak langsung.

Perbedaan dengan kerangka rujukan tetap ada tetapi alasannya bergeser. Kerangka itu mengandaikan servo
bertorsi dengan enkoder yang terbaca host. Di sini enkoder ada, tetapi \textbf{tertutup di dalam driver}:
host hanya mengirim pulsa langkah dan arah, tidak pernah membaca $\bm{q}$, dan tidak dapat memerintahkan
torsi sesaat. Karena itu \emph{Computed Torque Control} dan \emph{Sliding Mode Control} tetap tidak dapat
diterapkan apa adanya, kini karena antarmukanya, bukan karena ketiadaan sensor. Penggantinya adalah umpan
maju astronomis, pembentuk perintah \emph{input shaper} ZVD, dan tapis Kalman diperluas yang menaksir
simpangan di \emph{sisi beban}, yaitu bagian yang justru tidak terlihat oleh enkoder rotor.

Hasil utama: reduksi gigi mengubah anggaran torsi sumbu kemiringan dari $144\%$ (macet) menjadi $24\%$, dan
memperhalus microstep di beban dari $0{,}1125\degg$ menjadi $0{,}0225\degg$. Karena reduksi menaikkan pegas
magnet yang dirujuk ke beban melewati kekakuan poros, frekuensi alami kemiringan hanya bergeser $4{,}3\%$
terhadap seluruh rentang penguatan driver yang belum diketahui, sehingga datasheet yang masih ditunggu
praktis hanya memengaruhi azimuth. Tapis Kalman dengan LDR terkalibrasi menurunkan galat pointing dari
$2{,}52\degg$ menjadi $0{,}18\degg$; beban angin terburuk sebelas tahun memakai $9{,}2\%$ torsi tahan.
Panel diparkir datar saat hujan, dan itu terbukti tidak merugikan: pada mendung tebal panel datar justru
memanen $102{,}3\%$ dari panel yang menjejak.
\end{abstract}

\tableofcontents
\newpage

%=====================================================================
\section{Kerangka Koordinat dan Notasi}

\model{Kerangka toposentrik lokal North-West-Up (NWU)}

Kerangka rujukan memakai ENU; dokumen ini memakai NWU karena azimuth kompas kemudian muncul langsung
sebagai rotasi negatif terhadap sumbu tegak, tanpa konversi tambahan:
\begin{equation}
\hat{\bm{x}} = \text{Utara}, \qquad \hat{\bm{y}} = \text{Barat}, \qquad \hat{\bm{z}} = \text{atas (zenith)}.
\end{equation}

Sumbu engsel kemiringan \textbf{searah rel panjang}, yaitu sumbu $x$ alat, sehingga panel berguling ke kanan
dan kiri, bukan jungkat-jungkit sepanjang rel.

\begin{longtable}{p{3.6cm} p{1.6cm} p{9.6cm}}
\toprule
\textbf{Besaran} & \textbf{Lambang} & \textbf{Definisi} \\
\midrule
\endhead
Sudut sendi azimuth & $q_1$ & putar terhadap $\hat{\bm{z}}$, positif searah jarum jam dilihat dari atas, tanpa batas \\
Sudut sendi kemiringan & $q_2$ & $0$ = panel datar; batas $\pm 45\degg$ (ASUMSI) \\
Arah hadap rel & $\psi$ & $\psi = 90\degg$ (rel menghadap Utara, ASUMSI) \\
Arah miring kompas & $q_1 + \psi$ & arah panel condong, diukur searah jarum jam dari Utara \\
Azimuth matahari & $A$ & searah jarum jam dari Utara \\
Elevasi matahari & $e$ & dari horizon; zenith $z = 90\degg - e$ \\
\bottomrule
\end{longtable}

%=====================================================================
\section{Posisi Matahari dan Irradiance}

\subsection{Vektor arah matahari}

\model{NREL Solar Position Algorithm (SPA), Reda dan Andreas}

Kerangka rujukan memakai rumus deklinasi sederhana
$\sin\alpha_s = \sin\phi\sin\delta + \cos\phi\cos\delta\cos\omega$. Sistem ini memakai SPA melalui pustaka
\texttt{pvlib}, yang memperhitungkan nutasi, aberasi, dan refraksi atmosfer. Vektor satuan menuju matahari:
\begin{equation}
\shat = \begin{bmatrix} \cos e \cos A \\ -\cos e \sin A \\ \sin e \end{bmatrix}.
\end{equation}

\textbf{Tervalidasi:} selisih terhadap \texttt{astropy}, SolTrack, dan Pysolar \textbf{di bawah
$0{,}011\degg$}.

\subsection{Irradiance langit cerah dan tutupan awan}

\model{Langit cerah Ineichen-Perez}; \model{Faktor awan Kasten-Czeplak}; \model{Dekomposisi Erbs}

Irradiance horizontal global pada tutupan awan $N$ okta diperoleh dari nilai langit cerah:
\begin{equation}
G_h = G_{h,\mathrm{cs}} \left[ 1 - 0{,}75 \left( \frac{N}{8} \right)^{3{,}4} \right],
\qquad N \in [0,8].
\end{equation}
Pemisahan komponen langsung (DNI, $I_b$) dan baur (DHI, $I_d$) memakai korelasi Erbs melalui indeks
kecerahan $k_t = G_h / (I_{0}\cos z)$.

\subsection{Irradiance pada bidang panel}

\model{Transposisi isotropik Liu-Jordan}; \model{Koreksi sudut datang ASHRAE (Souka-Safwat)}

Sudut datang $\theta_i$ memenuhi $\cos\theta_i = \shat \cdot \nhat$. Faktor IAM yang memodelkan pantulan kaca:
\begin{equation}
K_\tau(\theta_i) = 1 - b_0 \left( \frac{1}{\cos\theta_i} - 1 \right), \qquad b_0 = 0{,}05 \ \text{(ASUMSI)}.
\end{equation}
Irradiance total pada bidang panel (\emph{plane of array}, POA):
\begin{equation}
G_{\mathrm{POA}} = \underbrace{I_b \cos\theta_i\, K_\tau}_{\text{langsung}}
+ \underbrace{I_d \frac{1 + \cos\beta}{2}}_{\text{baur langit}}
+ \underbrace{G_h\, \rho_g \frac{1 - \cos\beta}{2}}_{\text{pantulan tanah}},
\label{eq:poa}
\end{equation}
dengan $\beta$ kemiringan panel dari horizontal ($\cos\beta = n_z$) dan $\rho_g$ albedo tanah
($0{,}25$ ASUMSI).

\textbf{Catatan penting.} Dua suku terakhir Persamaan~\eqref{eq:poa} menarik titik optimum \emph{menjauh}
dari arah matahari, dan ke arah yang \emph{berlawanan} satu sama lain. Konsekuensinya dibahas di
Bagian~\ref{sec:arus}.

%=====================================================================
\section{Kinematika}

\subsection{Vektor normal panel}

\model{Kinematika rotasi benda tegar, rantai dua sendi putar}

Normal panel diperoleh dengan memutar vektor zenith $\bm{n}_0 = [0,0,1]^T$ berturut-turut terhadap sumbu
engsel kemiringan lalu sumbu azimuth:
\begin{equation}
\nhat = R_z(-q_1)\, R_x(q_2)\, \bm{n}_0
= \begin{bmatrix}
\sin q_2 \cos(q_1 + \psi) \\
-\sin q_2 \sin(q_1 + \psi) \\
\cos q_2
\end{bmatrix}.
\label{eq:fk}
\end{equation}

Basis panel yang menyertainya, dengan $a = q_1 + \psi$ dan $b = q_2$:
\begin{equation}
\xp = \begin{bmatrix} \cos b \cos a \\ -\cos b \sin a \\ -\sin b \end{bmatrix}
\quad (\text{arah \emph{chord}, melintang engsel}), \qquad
\yp = \begin{bmatrix} \sin a \\ \cos a \\ 0 \end{bmatrix}
\quad (\text{arah \emph{span}, sejajar engsel}).
\end{equation}

\textbf{Tervalidasi:} normal panel dari TF Gazebo berbanding Persamaan~\eqref{eq:fk} pada
$(q_1,q_2) = (0,30\degg)$, $(90\degg,30\degg)$, $(0,-30\degg)$ menunjukkan selisih $0{,}14$ sampai
$0{,}34\degg$, seluruhnya terjelaskan oleh defleksi gravitasi.

\subsection{Kinematika balik}

Untuk matahari pada $(A, z)$ terdapat \emph{dua} solusi,
\begin{equation}
(q_1, q_2) = (A - \psi,\ z) \qquad \text{atau} \qquad (A - \psi + 180\degg,\ -z),
\end{equation}
dan dipilih cabang yang kontinu terhadap langkah sebelumnya, sehingga melewati zenith tidak menuntut putaran
azimuth $180\degg$.

\subsection{Jacobian dan singularitas \emph{keyhole}}

\model{Jacobian analitik dan dekomposisi nilai singular}

Turunan Persamaan~\eqref{eq:fk} terhadap sudut sendi:
\begin{equation}
\frac{\partial \nhat}{\partial q_1} = -\sin q_2\, \yp, \qquad
\frac{\partial \nhat}{\partial q_2} = \xp,
\end{equation}
sehingga matriks Jacobian $\bm{J} = [\,\partial\nhat/\partial q_1 \ \ \partial\nhat/\partial q_2\,]
\in \mathbb{R}^{3\times 2}$ memiliki nilai singular $\sigma_1 = 1$ dan $\sigma_2 = |\sin q_2|$, dengan
bilangan kondisi
\begin{equation}
\kappa(\bm{J}) = \frac{1}{|\sin q_2|}.
\end{equation}

\begin{longtable}{c c c c c}
\toprule
$q_2$ & $1\degg$ & $5\degg$ & $30\degg$ & $45\degg$ \\
\midrule
$\kappa(\bm{J})$ & $57{,}3$ & $11{,}5$ & $2{,}0$ & $1{,}4$ \\
\bottomrule
\end{longtable}

Inilah \textbf{efek keyhole}, dinyatakan sebagai hilangnya rank, bukan sebagai gejala terpisah. Dampaknya
muncul kembali sebagai bobot sensor (Bagian~\ref{sec:fusi}) dan sebagai kovarians tapis
(Bagian~\ref{sec:ekf}), tanpa memerlukan aturan ambang buatan.

\subsection{Galat pointing dan rugi kosinus}

\model{Hukum kosinus Lambert}
\begin{equation}
\theta_{\mathrm{err}} = \arccos\!\left( \nhat \cdot \shat \right), \qquad
\text{rugi relatif komponen langsung} = 1 - \cos\theta_{\mathrm{err}}.
\label{eq:cosloss}
\end{equation}
Untuk ayunan sinusoidal beramplitudo $A_0$, rugi rata-rata satu siklus mendekati $A_0^2/4$ pada sudut kecil.

%=====================================================================
\section{Dinamika Multi-Benda}

\subsection{Persamaan gerak}

\model{Persamaan Euler-Lagrange untuk manipulator dua derajat kebebasan}

Dengan Lagrangian $\mathcal{L} = T - V$ dan koordinat umum $\bm{q} = [q_1, q_2]^T$, persamaan
$\frac{d}{dt}\frac{\partial\mathcal{L}}{\partial\dot{q}_i} - \frac{\partial\mathcal{L}}{\partial q_i}
= \tau_i$ menghasilkan bentuk baku manipulator:
\begin{equation}
\boxed{\;
\bm{M}(\bm{q})\,\ddot{\bm{q}} + \bm{C}(\bm{q},\dot{\bm{q}})\,\dot{\bm{q}} + \bm{g}(\bm{q})
+ \bm{F}(\dot{\bm{q}}) + \bm{\tau}_w = \bm{\tau} \; }
\label{eq:manip}
\end{equation}

\subsection{Matriks massa}

\model{Teorema sumbu sejajar (Steiner) atas tensor inersia CAD}

Inersia efektif sumbu azimuth bergantung kemiringan:
\begin{equation}
\bm{M}(\bm{q}) =
\begin{bmatrix}
J_1(q_2) & 0 \\ 0 & J_2
\end{bmatrix},
\qquad
J_1(q_2) = I_{z1} + \left( m_2 r_c^2 + I_{y2} \right)\cos^2 q_2 + I_{z2}\sin^2 q_2 ,
\qquad
J_2 = m_2 r_c^2 + I_{x2}.
\end{equation}
Bentuk ini dapat ditulis ulang sebagai $J_1(q_2) = a + b\cos^2 q_2$. Dari tensor inersia mesh CAD desain
baru, dengan panel $2{,}55$~kg menurut pelat nama:
\begin{equation}
J_1(q_2) = 0{,}1106 + 0{,}0408\cos^2 q_2 \ \mathrm{kg\,m^2}, \qquad
J_2 = 0{,}0587 \ \mathrm{kg\,m^2}, \qquad
\Delta I \equiv m_2 r_c^2 + I_{y2} - I_{z2} = 0{,}0408 \ \mathrm{kg\,m^2}.
\end{equation}
Nilai acuan: $J_1(0) = 0{,}1514$ dan $J_1(45\degg) = 0{,}1310 \ \mathrm{kg\,m^2}$. Massa yang berputar pada
sumbu kemiringan naik menjadi $3{,}782$~kg (dulu $3{,}055$) karena rangka desain baru membawa rel, braket
dan empat rakitan LDR; karena itu $J_2$ naik dari $0{,}0543$ ke $0{,}0587 \ \mathrm{kg\,m^2}$.

\textbf{Catatan reduksi gigi.} Nilai di atas adalah inersia \emph{di sisi beban}. Inersia rotor yang
dirujuk ke beban naik $N^2$, dan untuk sumbu kemiringan itu $1{,}1\times10^{-4} \times 25 = 2{,}75\times
10^{-3} \ \mathrm{kg\,m^2}$, masih dua orde di bawah $J_2$, jadi tidak mengubah mode lambat.

\subsection{Suku Coriolis dan sentripetal}

\model{Simbol Christoffel jenis pertama}

\begin{equation}
\bm{C}(\bm{q},\dot{\bm{q}}) =
\begin{bmatrix}
-\Delta I \sin(2q_2)\,\dot{q}_2 & 0 \\[2pt]
\tfrac{1}{2}\Delta I \sin(2q_2)\,\dot{q}_1 & 0
\end{bmatrix},
\end{equation}
sehingga $\dot{\bm{M}} - 2\bm{C}$ bersifat anti-simetris dan sistem bersifat pasif.

\textbf{Status: belum diterapkan dalam simulasi.} Besarnya sudah dihitung:

\begin{longtable}{p{6.4cm} r r}
\toprule
\textbf{Keadaan} & \textbf{Coriolis azimuth} & \textbf{\% torsi tahan} \\
\midrule
Jejak harian khas ($\dot q_1 = 0{,}08\degg$/s) & $4{,}3\times10^{-8}$ N\,m & $0{,}0000\%$ \\
Jejak laju maksimum & $5{,}4\times10^{-6}$ N\,m & $0{,}0002\%$ \\
Slew $30\degg$/s & $9{,}7\times10^{-3}$ N\,m & $0{,}32\%$ \\
\textbf{Slew $90\degg$/s} & $\bm{8{,}7\times10^{-2}}$ \textbf{N\,m} & $\bm{2{,}9\%}$ \\
\bottomrule
\end{longtable}

Tidak berpengaruh saat menjejak, tetapi mencapai $3{,}4\%$ saat slew cepat pada $q_2 = 45\degg$.

\subsection{Gravitasi}

\model{Momen gravitasi benda tegar tak seimbang}

Massa berputar pada sumbu kemiringan $m_2 = 3{,}782$~kg, dengan pusat massa bergeser hanya
$c_y = 0{,}23$~mm dan $c_z = 32{,}82$~mm dari sumbu engsel. \textbf{Desain baru jauh lebih seimbang}:
komponen $c_y$, yang menentukan momen pada posisi datar, turun dari $9{,}04$~mm menjadi $0{,}23$~mm.
\begin{equation}
\bm{g}(\bm{q}) = \begin{bmatrix} 0 \\ -m_2 g \left( c_y \cos q_2 - c_z \sin q_2 \right) \end{bmatrix},
\qquad g = 10 \ \mathrm{m/s^2} \ \text{(keputusan)}.
\label{eq:grav}
\end{equation}

\begin{longtable}{l r r}
\toprule
\textbf{Posisi} & \textbf{Momen gravitasi} & \textbf{Defleksi statis} \\
\midrule
$q_2 = 0$ (parkir) & $\bm{0{,}0088}$ \textbf{N\,m} & $0{,}0031\degg$; terukur Gazebo di bawah $0{,}001\degg$ \\
$q_2 = -45\degg$ & $\bm{0{,}884}$ \textbf{N\,m} (maksimum) & $0{,}315\degg$ \\
\bottomrule
\end{longtable}

\subsection{Gesekan}

\model{Gesekan viskos dan Coulomb (model Coulomb-viskos)}
\begin{equation}
\bm{F}(\dot{\bm{q}}) =
\begin{bmatrix} b_1 \dot q_1 + \tau_{c1}\operatorname{sgn}(\dot q_1) \\
b_2 \dot q_2 + \tau_{c2}\operatorname{sgn}(\dot q_2) \end{bmatrix},
\qquad b = 0{,}2 \ \mathrm{N\,m\,s/rad}, \quad \tau_c = 0{,}05 \ \mathrm{N\,m}.
\end{equation}
Keduanya \textbf{ASUMSI} dan belum diukur. Keduanya paling menentukan lama ayunan, dan dapat diperoleh
sekaligus lewat uji peluruhan bebas (Bagian~\ref{sec:asumsi}).

\subsection{Ruang keadaan tak linier}

\model{Representasi ruang keadaan orde satu}

Dengan $\bm{x} = [q_1, q_2, \dot q_1, \dot q_2]^T$:
\begin{equation}
\dot{\bm{x}} =
\begin{bmatrix}
\dot{\bm{q}} \\
\bm{M}^{-1}(\bm{q})\left[ \bm{\tau} - \bm{C}(\bm{q},\dot{\bm{q}})\dot{\bm{q}}
- \bm{g}(\bm{q}) - \bm{F}(\dot{\bm{q}}) - \bm{\tau}_w \right]
\end{bmatrix}.
\end{equation}

\subsection{Nilai eigen dan pembuktian kestabilan}

\model{Linearisasi Jacobian dan analisis nilai eigen sistem orde dua}

Di sekitar titik kerja, tiap sumbu menjadi sistem orde dua dengan kekakuan efektif $K$ (Bagian~\ref{sec:step})
dan redaman $b$. Dengan $\bm{x} = [q,\dot q]^T$:
\begin{equation}
\bm{F}_{\mathrm{dyn}} = \frac{\partial \bm{f}}{\partial \bm{x}} =
\begin{bmatrix} 0 & 1 \\ -K/J & -b/J \end{bmatrix},
\qquad
\lambda_{1,2} = -\zeta\omega_n \pm j\,\omega_n\sqrt{1-\zeta^2},
\end{equation}
dengan $\omega_n = \sqrt{K/J}$ dan $\zeta = b/(2\sqrt{KJ})$.

\begin{longtable}{l c c c c c}
\toprule
\textbf{Sumbu} & \textbf{Nilai eigen} & $\mathrm{Re}(\lambda)$ & $f_n$ & $\zeta$ & $\tau_{\text{peluruhan}}$ \\
\midrule
Azimuth & $-0{,}661 \pm 44{,}79j$ & $-0{,}661$ & $\bm{7{,}13}$ \textbf{Hz} & $0{,}0147$ & $1{,}51$ s \\
Kemiringan & $-1{,}704 \pm 52{,}34j$ & $-1{,}704$ & $\bm{8{,}33}$ \textbf{Hz} & $0{,}0325$ & $0{,}59$ s \\
\bottomrule
\end{longtable}

\textbf{Bagian nyata negatif untuk kedua sumbu, sehingga sistem stabil asimtotik dan panel pasti kembali ke
titik setimbang setelah bergerak.} Itu pembuktian formal yang diminta.

\model{Lewatan sistem orde dua kurang teredam}
\begin{equation}
M_p = \exp\!\left( \frac{-\pi\zeta}{\sqrt{1-\zeta^2}} \right)
\;\Longrightarrow\;
M_p^{\mathrm{az}} = 95{,}5\%, \qquad M_p^{\mathrm{el}} = 90{,}3\%.
\end{equation}
\textbf{Belum divalidasi ulang untuk konfigurasi ini.} Angka Gazebo yang pernah diukur ($90\%$ dan $82\%$
untuk langkah $1\degg$) berasal dari desain LAMA: penggerak langsung, $T_h = 3$~N\,m, loop terbuka. Dengan
reduksi gigi, motor $1{,}5$~N\,m dan loop tertutup, uji langkah dan uji peredam ZVD \emph{belum dijalankan
ulang}, sehingga tabel di atas adalah prediksi analitik tanpa jangkar pengukuran.

%=====================================================================
\section{Aktuator: Stepper Loop Tertutup dengan Reduksi Gigi}
\label{sec:step}

Bagian ini tidak ada dalam kerangka rujukan, tetapi menentukan seluruh rancangan kontrol di sini.

\subsection{Hukum torsi-sudut}

\model{Model torsi-sudut sinusoidal motor stepper hibrid}

Motor digerakkan pulsa posisi, bukan perintah torsi. Torsi pemulih bergantung selisih sudut listrik,
dan pada driver loop tertutup pegas magnet diperkuat penguatan posisi driver $K_p$:
\begin{equation}
\tau = \mathrm{Sat}\!\Big( T_h \sin( N_r e ) + K_p e, \ \pm T_h \Big),
\qquad e = \theta_{\mathrm{cmd}} - \theta_{\mathrm{rotor}},
\qquad T_h = 1{,}5 \ \mathrm{N\,m}, \quad N_r = 50,
\label{eq:stepper}
\end{equation}
dengan $K_p = 0$ mengembalikan hukum loop terbuka yang lama persis.

\textbf{Batas fisis yang menentukan $K_p$.} Torsi driver tetap dibatasi $T_h$, jadi kekakuan kecil-sinyal
hanya terwujud selama $|e| < T_h/(T_h N_r + K_p)$. Pada $K_p = 300$ jangkauan linier itu $0{,}229\degg$,
sekitar dua microstep. Limit ``rotor terpaku'' ($K_p \to \infty$) \textbf{tidak fisis}: jangkauan liniernya
nol, sehingga kekakuan yang dijanjikannya tidak pernah terwujud pada amplitudo mana pun.

\subsection{Reduksi gigi}

\model{Rantai roda gigi lurus satu tingkat}

Sumbu kemiringan memakai satu tingkat reduksi; azimuth tetap penggerak langsung.
\begin{equation}
N = \frac{z_2}{z_1} = \frac{30}{6} = 5{,}0, \qquad m = 2{,}75 \ \mathrm{mm},
\qquad d_1 = 16{,}5, \quad d_2 = 82{,}5 \ \mathrm{mm}.
\end{equation}
Rasio ini \textbf{diukur langsung di alat} (lima putaran pinion untuk satu putaran gigi besar), dan
diameter luar terukur $22$ dan $88$~mm cocok dengan $OD = m(z+2)$ pada modul $2{,}75$. Jarak sumbu
$49{,}5$~mm berbanding $50{,}2$~mm pada CAD.

Reduksi melipatkan torsi $N$ kali dan membagi langkah $N$ kali:
\begin{equation}
T_{h,L} = N T_h = 7{,}5 \ \mathrm{N\,m}, \qquad
\Delta_{\mu,L} = \frac{\Delta_\mu}{N} = 0{,}0225\degg .
\end{equation}

\subsection{Kekakuan seri}

\model{Pegas seri (pegas magnet dirujuk ke beban, seri dengan poros)}

Pegas magnet yang dirujuk ke sisi beban naik $N^2$; poros sudah berada di sisi beban sehingga tidak
diskalakan:
\begin{equation}
K = \left( \frac{1}{(T_h N_r + K_p) N^2} + \frac{1}{k} \right)^{-1}
\;\Longrightarrow\;
K_{\mathrm{az}} = 303{,}80, \qquad K_{\mathrm{el}} = 160{,}99 \ \mathrm{N\,m/rad},
\end{equation}
dengan $k_{\mathrm{az}} = 1600$ dan $k_{\mathrm{el}} = 163{,}8$~N\,m/rad. Nilai poros kemiringan itu
dihitung dari panjang \textbf{terukur} di CAD, $194$~mm untuk poros baja $8$~mm, lewat $k = GJ/L$.

\textbf{Akibat yang tidak intuitif dan penting untuk tesis ini.} Pada sumbu bergigi, $(T_h N_r + K_p)N^2
= 1875$~N\,m/rad jauh melewati poros $163{,}8$, sehingga \textbf{poros yang menentukan} dan penguatan
driver nyaris tidak berpengaruh. Frekuensi alami kemiringan hanya bergerak dari $8{,}06$ ke $8{,}41$~Hz
untuk $K_p$ dari nol sampai tak hingga, yaitu sebaran $4{,}3\%$. Pada azimuth yang penggerak langsung,
rentang yang sama memberi $3{,}46$ sampai $16{,}36$~Hz, sebaran $373\%$. Jadi datasheet driver yang masih
ditunggu praktis hanya memengaruhi azimuth.

\subsection{Kuantisasi microstep dan kejenuhan}

\model{Kuantisasi seragam (zero-order hold spasial)}
\begin{equation}
\theta_{\mathrm{cmd}} = \Delta_\mu \left\lfloor \frac{\theta_{\mathrm{ref}}}{\Delta_\mu} \right\rceil,
\qquad \Delta_\mu = \frac{1{,}8\degg}{16} = 0{,}1125\degg .
\end{equation}
Tangga inilah \textbf{sumber eksitasi ayunan saat menjejak}, bukan gangguan luar.

Pada sumbu bergigi, kuantisasi yang terlihat beban menjadi $\Delta_{\mu,L} = 0{,}0225\degg$, lima kali
lebih halus.

Bila torsi luar melewati $T_{h,L}$ motor kehilangan sinkron. \textbf{Berbeda dengan desain lama, kini
driver mengoreksi sendiri} karena enkoder rotor melihat selisihnya; yang tersisa adalah \emph{alarm galat
ikut} ketika selisih melewati ambang, bukan hanyutan permanen yang diam-diam menumpuk.

\textbf{Batas model.} Hukum Persamaan~\eqref{eq:stepper} hanya bernilai tunggal bila $k > T_{h,L} N_{r,L}$.
Dengan reduksi, ruas kanan menjadi $2067$~N\,m/rad melawan poros $163{,}8$, sehingga syarat itu
\textbf{gagal} dan kurva torsi kuasi-statis menjadi bernilai ganda. Karena itu simulasi desain ini memakai
model aktuator linier (kekakuan $K$ dengan kejenuhan $T_{h,L}$), bukan hukum sinus loop terbuka.

%=====================================================================
\section{Strategi Kontrol}

Kerangka rujukan mengusulkan \emph{Computed Torque Control},
$\bm{\tau} = \bm{M}(\ddot{\bm{q}}_d + K_v\dot{\bm{e}} + K_p\bm{e}) + \bm{C}\dot{\bm{q}} + \bm{g} + \bm{F}$,
dan \emph{Sliding Mode Control} dengan permukaan $S = \dot{\bm{e}} + \Lambda\bm{e}$. Keduanya menuntut
$\bm{q}$ dan $\dot{\bm{q}}$ terukur tiap siklus serta kemampuan mengatur torsi sesaat.

\textbf{Alat ini kini punya enkoder, tetapi tetap tidak memenuhi syarat itu, dan alasannya bergeser.}
Enkoder berada \emph{di dalam} driver loop tertutup: host mengirim pulsa langkah dan arah, tidak pernah
membaca $\bm{q}$, dan tidak dapat memerintahkan torsi sesaat. Lebih jauh, enkoder mengukur \emph{rotor},
bukan beban, sehingga berada di sisi yang salah dari kelenturan poros dan reduksi gigi. Jadi premis
``tidak ada pengukuran'' kini hanya separuh benar, sedangkan premis ``tidak ada antarmuka torsi'' tetap
berlaku penuh, dan premis kedua itu saja sudah cukup untuk menyingkirkan CTC maupun SMC.

Penggantinya tiga lapis berikut.

\subsection{Umpan maju astronomis}

Lintasan sendi dihitung dari SPA lewat kinematika balik, dibatasi $|q_2| \le 45\degg$.

\subsection{Profil gerak}

\model{Profil kecepatan trapesium dengan batas percepatan}
\begin{equation}
\dot q_{\mathrm{des}} = \operatorname{sgn}(e)\min\!\left( v_{\max}, \sqrt{2 a_{\max} |e|} \right),
\qquad v_{\max} = 90\degg/\mathrm{s}, \quad a_{\max} = 300\degg/\mathrm{s^2},
\end{equation}
dengan $e = q_{\mathrm{ref}} - q$. Bentuk akar memastikan berhenti tepat di sasaran tanpa lewatan perintah.

\subsection{Peredam ayunan}

\model{Input shaper ZV dan ZVD (Singer dan Seering)}

Deret impuls dikonvolusikan dengan perintah agar getaran sisa saling meniadakan. Dengan
$K_s = \exp\!\big(-\pi\zeta/\sqrt{1-\zeta^2}\big)$ dan jeda $\Delta t = T_d/2 = \big(2 f_n\sqrt{1-\zeta^2}\big)^{-1}$:
\begin{align}
\text{ZV:}\quad & \bm{A} = \frac{1}{1+K_s}\begin{bmatrix} 1 & K_s \end{bmatrix},
&& \bm{t} = \begin{bmatrix} 0 & \Delta t \end{bmatrix}, \\
\text{ZVD:}\quad & \bm{A} = \frac{1}{(1+K_s)^2}\begin{bmatrix} 1 & 2K_s & K_s^2 \end{bmatrix},
&& \bm{t} = \begin{bmatrix} 0 & \Delta t & 2\Delta t \end{bmatrix}.
\end{align}
ZVD menambah satu impuls agar tahan terhadap galat taksiran $f_n$, dengan ongkos satu $\Delta t$ tambahan.

\textbf{Hasil terukur di Gazebo}, lompatan azimuth $181\degg$, lima ulangan:

\begin{longtable}{l r r}
\toprule
\textbf{Strategi} & \textbf{Lewatan median} & \textbf{Lewatan terburuk} \\
\midrule
Hanya profil trapesium & $0{,}58\%$ ($1{,}05\degg$) & $\bm{10{,}71\%}$ ($\bm{19{,}35\degg}$) \\
\textbf{Profil + ZVD} & $0{,}11\%$ ($0{,}20\degg$) & $\bm{1{,}27\%}$ ($\bm{2{,}30\degg}$) \\
\bottomrule
\end{longtable}

Perbandingan langkah tunggal (model Python / Gazebo):

\begin{longtable}{l r r r r}
\toprule
\textbf{Uji} & \textbf{Mentah} & \textbf{Trapesium} & \textbf{ZV} & \textbf{ZVD} \\
\midrule
Azimuth $90\degg$ & $52{,}43 / 52{,}41\%$ & $0{,}88 / 0{,}77\%$ & $3{,}07 / 0{,}00\%$ & $0{,}26 / 0{,}00\%$ \\
Azimuth $5\degg$ & $82{,}60 / 82{,}78\%$ & $15{,}46 / 15{,}77\%$ & $0{,}32 / 0{,}40\%$ & $0{,}97 / 0{,}85\%$ \\
\bottomrule
\end{longtable}

\textbf{Batas penting.} Input shaper \emph{tidak berguna saat menjejak pelan}: satu microstep tidak dapat
dipecah menjadi dua setengah langkah, sehingga setelah pembulatan perintahnya identik. Satu-satunya tuas
saat menjejak adalah memperhalus $\Delta_\mu$.

Sebab tambahan: pada lompatan besar pengendali \emph{jenuh} di $T_h$ sehingga geraknya bang-bang, dan
input shaper yang mengandaikan sistem linier sama sekali tidak berefek (terukur: lewatan $71{,}2\%$ identik
dengan dan tanpa ZV). Karena itu shaper harus dipasang \emph{di atas} profil trapesium.

%=====================================================================
\section{Model Listrik Panel Surya}
\label{sec:arus}

\subsection{Persamaan dioda tunggal}

\model{Model dioda tunggal lima parameter (De Soto, Klein, Beckman); persamaan Shockley}

\begin{equation}
I = I_L - I_0\left[ \exp\!\left( \frac{V + I R_s}{a} \right) - 1 \right] - \frac{V + I R_s}{R_{sh}},
\qquad a = \frac{n N_s k T_c}{q},
\label{eq:diode}
\end{equation}
dengan $I_L$ arus cahaya, $I_0$ arus jenuh balik, $R_s$ tahanan seri, $R_{sh}$ tahanan bocor, $n$ faktor
ideal, dan $N_s$ jumlah sel seri.

Parameter dicocokkan ke pelat nama MOS C-50-18-M dengan empat syarat: titik hubung singkat, titik rangkaian
terbuka, titik daya maksimum, dan $\partial P/\partial V = 0$ di titik daya maksimum.

\begin{longtable}{l r l r}
\toprule
\textbf{Parameter} & \textbf{Nilai} & \textbf{Reproduksi STC} & \textbf{Model} \\
\midrule
$I_L$ & $3{,}0092$ A & $I_{sc}$ & $3{,}000$ A \\
$I_0$ & $7{,}02\times10^{-10}$ A & $V_{oc}$ & $22{,}50$ V \\
$R_s$ & $0{,}372\ \Omega$ & $I_{mp}$ & $2{,}700$ A \\
$R_{sh}$ & $121{,}5\ \Omega$ & $V_{mp}$ & $18{,}50$ V \\
$n$ & $1{,}1$ & $P_{\max}$ & $49{,}95$ W \\
\bottomrule
\end{longtable}

Faktor ideal $n$ tidak dapat ditentukan dari pelat nama (empat persamaan, lima anu); diperiksa bahwa
$n \in [1{,}0;\,1{,}5]$ mengubah arus hanya $0{,}13\%$, sehingga pilihannya tidak kritis.

\subsection{Suhu sel}

\model{Model NOCT (Nominal Operating Cell Temperature)}
\begin{equation}
T_c = T_a + \frac{\mathrm{NOCT} - 20\degg\mathrm{C}}{800\ \mathrm{W/m^2}}\, G_{\mathrm{POA}},
\qquad \mathrm{NOCT} = 45\degg\mathrm{C} \ \text{(ASUMSI)}.
\end{equation}
Koefisien daya terhitung $-0{,}497\%/\degg$C, wajar untuk silikon. Efisiensi modul STC $16{,}75\%$.

\subsection{Arus sebagai sensor arah: dua batas mendasar}

\textbf{Pertama, memaksimalkan arus tidak mengarahkan panel ke matahari.} Akibat dua suku sebaran pada
Persamaan~\eqref{eq:poa}, puncak arus bergeser:

\begin{longtable}{l r l}
\toprule
\textbf{Albedo $\rho_g$} & \textbf{Pergeseran puncak} & \textbf{Contoh permukaan} \\
\midrule
$0{,}10$ & $-2{,}75\degg$ & aspal, tanah basah \\
$0{,}25$ & $-0{,}10\degg$ & nilai bawaan model \\
$0{,}40$ & $+2{,}75\degg$ & beton terang, pasir \\
\bottomrule
\end{longtable}

Dengan komponen langsung saja puncaknya tepat di $0\degg$. \textbf{Albedo lokasi harus diukur} bila arus
dipakai sebagai acuan posisi.

\textbf{Kedua, arus buta terhadap galat kecil.} Karena Persamaan~\eqref{eq:cosloss} bersifat kuadratik di
dekat puncak, simpangan $2\degg$ hanya mengubah arus $0{,}09\%$, jauh di bawah toleransi pabrik $\pm3\%$.
Lebih jauh, perubahan arus akibat awan \textbf{dua sampai tiga orde} lebih besar daripada akibat salah arah:

\begin{longtable}{l r r}
\toprule
\textbf{Tutupan awan} & \textbf{Turun karena awan} & \textbf{Turun karena salah arah $5\degg$} \\
\midrule
Okta 2 & $-1{,}0\%$ & $-0{,}29\%$ \\
Okta 4 & $-10{,}7\%$ & $-0{,}06\%$ \\
Okta 6 & $-40{,}8\%$ & $+0{,}96\%$ \\
\bottomrule
\end{longtable}

Pendakian bukit (\emph{perturb and observe}) memakai arus berhasil saat cerah dan tenang (berhenti
$0{,}30\degg$ dari puncak) tetapi \textbf{tersesat sampai $51\degg$} ketika awan menipis $2\%$ per detik.

%=====================================================================
\section{Sensor LDR dan Penggabungan Berbobot}
\label{sec:fusi}

\subsection{Model pembacaan}

\model{Pembagi tegangan dengan LDR hukum pangkat, dan sekat pembayang}

Hambatan LDR mengikuti hukum pangkat terhadap iluminansi, sehingga tegangan keluaran pembagi bersifat jenuh:
\begin{equation}
x = \left( \frac{E}{L_0} \right)^{\gamma}, \qquad
V = V_{cc}\,\frac{x}{1+x} + \nu, \qquad \nu \sim \mathcal{N}(0, \sigma_{\mathrm{ADC}}^2),
\end{equation}
dengan $\gamma = 0{,}7$ dan $L_0 = 20000$ lux (keduanya ASUMSI). Sekat membayangi komponen \emph{langsung}
sebanding $r\tan\varepsilon$ dengan $\varepsilon$ sudut miring matahari terhadap sisi LDR.

Sinyal selisih yang dipakai:
\begin{equation}
e_x = \frac{V_{x+} - V_{x-}}{V_{x+} + V_{x-}}, \qquad
e_y = \frac{V_{y+} - V_{y-}}{V_{y+} + V_{y-}} .
\label{eq:ldr}
\end{equation}
Pembagian dengan jumlah membuat sinyal tahan terhadap perubahan irradiance menyeluruh.

\subsection{Pemisahan derau dan bias}

Ketidakpastian taksiran sudut adalah derau sinyal dibagi kepekaan, $\sigma_q = \sigma_e / |\partial e/\partial q|$.
Diukur Monte Carlo dengan ketidakcocokan antar sensor diundi \emph{sekali per perangkat}:

\begin{longtable}{l r r r r}
\toprule
\textbf{Langit} & \textbf{Derau acak $q_2$} & \textbf{Bias $q_2$} & \textbf{Derau acak $q_1$} & \textbf{Bias $q_1$} \\
\midrule
Cerah & $0{,}44\degg$ & $\bm{2{,}86\degg}$ & $0{,}60\degg$ & $\bm{4{,}68\degg}$ \\
Mendung okta 6 & $0{,}69\degg$ & $\bm{5{,}20\degg}$ & $0{,}94\degg$ & $\bm{8{,}54\degg}$ \\
\bottomrule
\end{longtable}

\textbf{Yang menguasai bukan derau, melainkan bias}, lima sampai sepuluh kali lebih besar. Derau dapat
ditekan dengan merata-rata; bias hanya hilang dengan kalibrasi.

\subsection{Penggabungan bobot kebalikan ragam}

\model{Penaksir bobot kebalikan ragam (inverse-variance weighting, BLUE)}

Untuk dua taksiran tak bias dan bebas:
\begin{equation}
\hat{x} = \frac{\sigma_2^2 \hat{x}_1 + \sigma_1^2 \hat{x}_2}{\sigma_1^2 + \sigma_2^2},
\qquad \frac{1}{\sigma^2} = \frac{1}{\sigma_1^2} + \frac{1}{\sigma_2^2}.
\end{equation}

\begin{longtable}{l r r r r}
\toprule
\textbf{Ketidakcocokan LDR} & $\sigma_{q_2}$ & $\sigma_{q_1}$ & \textbf{Bobot LDR $q_2$} & \textbf{Bobot LDR $q_1$} \\
\midrule
$3\%$ (tanpa kalibrasi) & $2{,}92\degg$ & $4{,}08\degg$ & $\bm{20{,}9\%}$ & $\bm{11{,}9\%}$ \\
$1\%$ (kalibrasi kasar) & $1{,}03\degg$ & $1{,}46\degg$ & $67{,}8\%$ & $51{,}4\%$ \\
$0{,}3\%$ (kalibrasi teliti) & $0{,}51\degg$ & $0{,}71\degg$ & $\bm{89{,}5\%}$ & $\bm{81{,}8\%}$ \\
\bottomrule
\end{longtable}

Bobot LDR untuk azimuth turun \emph{sendiri} dari $9{,}2\%$ pada $q_2 = 45\degg$ menjadi $0{,}5\%$ pada
$q_2 = 10\degg$, karena $\partial e_y/\partial q_1 \propto \sin q_2$. Efek keyhole muncul kembali sebagai
bobot, tanpa aturan ambang buatan.

\subsection{Prosedur kalibrasi}

\textbf{Prosedur A, cahaya baur merata.} Sekat hanya membayangi komponen langsung, sehingga di bawah cahaya
baur murni keempat LDR menerima iluminansi sama berapa pun arah panel. Faktor koreksi dicari di ruang
$x = V/(V_{cc}-V)$:
\begin{equation}
m_i = \left( \frac{x_i}{\bar{x}} \right)^{1/\gamma}, \qquad c_i = 1/m_i .
\end{equation}
Diperiksa: faktor koreksi identik sampai satu bagian per juta pada tiga pose berbeda. Di lapangan cukup
menunggu mendung tebal (okta 8 menyisakan DNI hanya $2{,}2\%$).

\textbf{Prosedur B, regresi sepanjang hari.} Salah pasang menghasilkan $e = \kappa\,\delta$ yang ikut
berubah ketika kepekaan $\kappa$ berubah, sedangkan ketidakcocokan menghasilkan offset hampir tetap:
\begin{equation}
e[k] = \kappa[k]\,\delta + c + \varepsilon[k].
\end{equation}
Kepekaan azimuth berubah $25{,}6$ kali sepanjang hari, kemiringan hanya $1{,}1$ kali. Hasilnya bias azimuth
terpisahkan dengan baik (taksiran $1{,}50\degg$ dari sebenarnya $1{,}50\degg$), bias kemiringan tidak
($0{,}84$ dari $1{,}00$, bilangan kondisi $23936$).

%=====================================================================
\section{Beban Angin dari Data Terukur}

\subsection{Sumber dan koreksi ketinggian}

\model{Reanalisis Open-Meteo (ERA5); profil pangkat Hellmann}

Data: $102{,}984$ jam reanalisis per jam untuk titik PENS, 2015 sampai 2026, lengkap dengan arah dan
hembusan. Reanalisis berlaku pada $10$~m; panel hanya $0{,}3$~m di atas tanah:
\begin{equation}
v(z) = v_{10}\left( \frac{z}{10} \right)^{\alpha}, \qquad \alpha = 0{,}30 \ \text{(ASUMSI, kampus)}.
\end{equation}

\begin{longtable}{l r r}
\toprule
& \textbf{Pada 10 m} & \textbf{Di panel 0,3 m} \\
\midrule
Kecepatan rata-rata & $2{,}83$ m/s & $0{,}99$ m/s \\
Kecepatan maksimum & $9{,}33$ m/s & $3{,}26$ m/s \\
\textbf{Hembusan maksimum 11 tahun} & $\bm{17{,}81}$ \textbf{m/s} & $\bm{6{,}22}$ \textbf{m/s} \\
\bottomrule
\end{longtable}

Arah didominasi Timur ($34{,}3\%$), pola monsun. Maksimum hembusan tahunan sangat stabil (rata $16{,}28$ m/s,
simpangan baku $1{,}29$), karena Surabaya terlalu dekat ekuator untuk siklon tropis.

\subsection{Model torsi berarah}

\model{Gaya normal pelat datar}

Angin dari bearing $\varphi$ berjalan menuju $\varphi + 180\degg$:
\begin{equation}
\what = \big[ -\cos\varphi,\ \sin\varphi,\ 0 \big]^T,
\qquad
F = \tfrac{1}{2}\rho v^2 A\, C_n \left( \what \cdot \nhat \right).
\end{equation}
Dengan pergeseran pusat tekanan $e_c$ searah \emph{chord} dan $e_s$ searah \emph{span}, dan karena $\bm{F}$
sejajar $\nhat$ sehingga $\bm{r}\times\bm{F} = -e_c F\,\yp + e_s F\,\xp$:
\begin{equation}
\boxed{\;
\tau_{\mathrm{el}} = -e_c F, \qquad
\tau_{\mathrm{az}} = -e_s F \sin q_2 \; }
\label{eq:wind}
\end{equation}

\textbf{Torsi azimuth nol saat panel datar}, karena gaya tegak lurus panel datar tidak dapat memuntir sumbu
tegak. Model buta arah pada kerangka rujukan tidak memiliki sifat ini.

\textbf{Tervalidasi di Gazebo:} pengali $1$, $2$, dan $3$ memberi $0{,}45\%$, $1{,}80\%$, dan $4{,}04\%$
torsi tahan; rasio $4$ dan $9$ terbaca persis sesuai $v^2$. Tanda torsi juga berbalik antara pagi dan siang
karena orientasi panel berubah.

\subsection{Margin terhadap lepas sinkron}

\begin{longtable}{l r r}
\toprule
\textbf{Kasus} & $v$ \textbf{di panel} & \textbf{\% torsi tahan} \\
\midrule
\textbf{Terburuk 11 tahun} & $6{,}22$ m/s & $\bm{9{,}2\%}$ \\
Pengali $1{,}5$ & $9{,}33$ m/s & $20{,}8\%$ \\
Pengali $2$ & $12{,}44$ m/s & $37{,}0\%$ \\
Pengali $3$ & $18{,}66$ m/s & $83{,}1\%$ \\
\bottomrule
\end{longtable}

\textbf{Angin bukan risiko bagi alat ini}, dengan margin sekitar sebelas kali. Preset sintetis
(\emph{tenang}, \emph{berangin}, \emph{badai}) karena itu dibuang dan diganti pengali terhadap data nyata.

%=====================================================================
\section{Hujan: Parkir Datar, Kaca Basah, dan Pendinginan}

\model{Reanalisis presipitasi Open-Meteo (ERA5)}

Sumber sama dengan angin: $102{,}984$ jam per jam untuk titik PENS, 2015 sampai 2026, kolom presipitasi,
hujan dan kelembapan. Curah $1792$~mm per tahun, dan pola monsun jelas: $43{,}7\%$ jam siang basah pada
Februari berbanding $2{,}2\%$ pada Agustus. Ambang kebijakan $0{,}2$~mm/jam dengan tahan $1$~jam setelah
hujan berhenti (ASUMSI, menunggu sensor hujan di alat).

\subsection{Mengapa parkir datar tidak merugikan}

Secara optik hujan adalah mendung tebal. Pada okta 8 komponen langsung runtuh ke sekitar $4$~W/m$^2$,
sehingga hampir seluruh cahaya menjadi baur. Panel datar melihat \emph{seluruh} kubah langit, panel miring
melihat lebih sedikit. Energi sehari pada bidang panel (7 Okt 2026, kWh/m$^2$):

\begin{longtable}{c r r r}
\toprule
\textbf{okta} & \textbf{menjejak} & \textbf{datar} & \textbf{datar/menjejak} \\
\midrule
$0$ (cerah) & $7{,}70$ & $7{,}25$ & $94{,}1\%$ \\
$4$ & $7{,}09$ & $6{,}73$ & $94{,}9\%$ \\
$6$ & $5{,}29$ & $5{,}20$ & $98{,}3\%$ \\
$\bm{8}$ \textbf{(hujan)} & $1{,}77$ & $1{,}81$ & $\bm{102{,}3\%}$ \\
\bottomrule
\end{longtable}

Jadi parkir datar saat hujan justru menang tipis, di samping membuat air mengalir keluar dan mengurangi
luas hadap terhadap angin. Hanya kemiringan yang diratakan; azimuth dibiarkan di tempat karena datar
berarti $q_2 = 0$ berapa pun azimuthnya.

\subsection{Kaca basah: lapisan air menolong, tetesan merugikan}

\model{Pantulan Fresnel pada antarmuka berlapis}

Ini sering disalahpahami. Lapisan air rata \textbf{mengurangi} pantulan, karena indeks bias air berada di
antara udara dan kaca sehingga bekerja sebagai lapisan pencocok indeks:
\begin{equation}
R(n_1,n_2) = \left(\frac{n_1-n_2}{n_1+n_2}\right)^2, \qquad
n_{\mathrm{udara}} = 1{,}0, \quad n_{\mathrm{air}} = 1{,}333, \quad n_{\mathrm{kaca}} = 1{,}52 .
\end{equation}
Pada datang tegak: kering memantulkan $4{,}26\%$, basah $2{,}46\%$, sehingga transmisi \emph{naik}
$1{,}88\%$. Suku ini diturunkan, bukan ditebak.

Yang merugikan adalah \textbf{tetesan} yang menghamburkan dan membayangi, dan itulah satu-satunya suku
empiris, $5\%$ (ASUMSI; pustaka $2$ sampai $8\%$). Faktor transmisi basah bersih:
\begin{equation}
T_{\mathrm{basah}} = \frac{(1-R_{\mathrm{ua}})(1-R_{\mathrm{ak}})}{1-R_{\mathrm{uk}}}\,(1-\epsilon_{\mathrm{tetes}})
= 0{,}9679 .
\end{equation}

\subsection{Pendinginan, dan jebakan hitung ganda}

Model NOCT \textbf{sudah} menurunkan suhu sel saat hujan semata karena irradiance jatuh: pada udara
$30\degg$C, POA $1000$ memberi $61{,}2\degg$C sedangkan POA $250$ memberi $37{,}8\degg$C. Pendinginan
\emph{tambahan} oleh air hanya sisa $37{,}8 \to 30\degg$C, yaitu sekitar $3{,}9\%$ daya pada
$-0{,}497\%/\degg$C. Mengutip belasan persen untuk pendinginan hujan berarti menghitung ganda efek
irradiance yang sudah masuk lewat okta.

\textbf{Hasil bersih pada irradiance sama} (okta 8): transmisi $0{,}9679$ dan suhu sel $37{,}8 \to
30{,}0\degg$C memberi $P_{mp}$ $11{,}265 \to 11{,}379$~W, yaitu $+1{,}01\%$. Panel basah menghasilkan
sedikit \emph{lebih} daripada panel kering pada irradiance sama. Hujan merugikan lewat runtuhnya
irradiance, bukan lewat basahnya.

\textbf{Tervalidasi di Gazebo} pada 18 Februari 2026 (hari siang terbasah tahun itu, $18{,}1$~mm), okta 8:
saat penanda hujan aktif, transmisi terbaca $0{,}9679$, suhu sel $30{,}0\degg$C, dan kemiringan parkir di
$0{,}00\degg$; saat kering, transmisi $1{,}0$, suhu sel $33$ sampai $36{,}5\degg$C, dan panel menjejak.

%=====================================================================
\section{Osilasi dan Dampaknya pada Energi}

Ayunan diukur pada laju penuh $1000$~Hz, dengan gerak jejak lambat dibuang lewat penapis lolos tinggi
(rata-rata bergerak di atas $2$~Hz):

\begin{longtable}{l r r}
\toprule
\textbf{Keadaan} & \textbf{Ayunan puncak ke puncak} & \textbf{Frekuensi terukur} \\
\midrule
Menjejak dekat transit zenith & az $1{,}90\degg$, el $0{,}72\degg$ & $5{,}19$ dan $6{,}57$ Hz \\
Setelah lompat jam besar & az $4{,}53\degg$, el $5{,}44\degg$ & \\
Sembilan detik kemudian & $0{,}000\degg$ & \\
\bottomrule
\end{longtable}

Frekuensi terukur $5{,}19$ dan $6{,}57$~Hz \textbf{cocok dengan prediksi nilai eigen} $4{,}79$ dan
$6{,}89$~Hz, dan peluruhan sembilan detik sesuai konstanta waktu $1{,}51$~s.

\begin{longtable}{l r r}
\toprule
\textbf{Penyebab rugi energi} & \textbf{Besarnya} & \\
\midrule
Suhu sel $55\degg$C dibanding $25\degg$C & $\bm{14{,}9\%}$ & \\
Toleransi pabrik panel & $3{,}0\%$ & \\
Defleksi gravitasi $0{,}72\degg$ & $0{,}008\%$ & \\
\textbf{Ayunan motor saat menjejak} & $\bm{0{,}007\%}$ & \\
\bottomrule
\end{longtable}

\textbf{Osilasi hampir tidak merugikan energi}, dua ribu kali lebih kecil daripada pengaruh suhu. Dampak
harian lompatan besar lima kali sehari hanya $0{,}19$~mWh dari panen $431$~Wh. Peredam ZVD tetap diperlukan,
tetapi alasannya beban mekanik, risiko lepas sinkron, dan derau ke sensor, bukan panen energi.

%=====================================================================
\section{Tapis Kalman Diperluas sebagai Penaksir Sisi Beban}
\label{sec:ekf}

\subsection{Keadaan yang ditaksir}

Kerangka rujukan mengusulkan UKF menaksir parameter fisis
$\bm{\theta} = [m_2, r_c, b_1, b_2, C_{m,A}, C_{m,E}]$. Enkoder driver sudah menutup loop di sisi rotor,
sehingga lepas langkah terkumpul bukan lagi masalah. Yang tersisa, dan justru tidak terlihat enkoder,
adalah segala sesuatu di \emph{seberang} gigi dan poros: puntiran poros, backlash gigi, salah pasang
mekanik, dan bias LDR. Karena itu yang ditaksir adalah
\emph{simpangan} antara perintah dan beban:
\begin{equation}
\bm{b} = \begin{bmatrix} b_{\mathrm{az}} \\ b_{\mathrm{el}} \end{bmatrix}
= \bm{q}_{\mathrm{beban}} - \bm{q}_{\mathrm{perintah}},
\end{equation}
yang memuat salah pasang mekanik, \textbf{lepas langkah terkumpul}, dan sisa defleksi.

\subsection{Persamaan tapis}

\model{Extended Kalman Filter, bentuk kovarians Joseph}

Model proses jalan acak ($\bm{F} = \bm{I}$) karena lepas langkah tidak dapat diramalkan:
\begin{align}
\text{Prediksi:} \quad & \hat{\bm{b}}^- = \hat{\bm{b}}, \qquad
\bm{P}^- = \bm{P} + \bm{Q}\,\Delta t, \\
\text{Inovasi:} \quad & \bm{y} = \bm{z} - h(\hat{\bm{b}}^-), \qquad
\bm{H} = \left.\frac{\partial h}{\partial \bm{b}}\right|_{\hat{\bm{b}}^-}, \qquad
\bm{S} = \bm{H}\bm{P}^-\bm{H}^T + \bm{R}, \\
\text{Penguatan:} \quad & \bm{K} = \bm{P}^-\bm{H}^T\bm{S}^{-1}, \qquad
\hat{\bm{b}} = \hat{\bm{b}}^- + \bm{K}\bm{y}, \\
\text{Kovarians:} \quad & \bm{P} = (\bm{I}-\bm{K}\bm{H})\bm{P}^-(\bm{I}-\bm{K}\bm{H})^T + \bm{K}\bm{R}\bm{K}^T .
\end{align}
Bentuk Joseph dipakai agar $\bm{P}$ tetap simetris dan definit positif.

Fungsi pengukuran $h$ adalah model LDR Persamaan~\eqref{eq:ldr}, sehingga $\bm{H}$ mengandung Jacobian
kinematik; itulah pertemuan antara singularitas keyhole dan kovarians tapis.

\model{Uji pencilan khi-kuadrat atas inovasi ternormalkan (NIS)}
\begin{equation}
\epsilon = \bm{y}^T \bm{S}^{-1} \bm{y} \;\sim\; \chi^2_{m}, \qquad
\text{pengukuran ditolak bila } \epsilon > \epsilon_{\max} = 25 .
\end{equation}

\subsection{Hasil terukur}

Hari 7 Oktober 2026, lepas langkah azimuth $7{,}2\degg$ disuntikkan di tengah hari; galat diukur hanya saat
matahari terjangkau:

\begin{longtable}{l r r}
\toprule
\textbf{Susunan} & \textbf{Galat rata-rata} & \textbf{p95} \\
\midrule
Astro saja, tanpa tapis & $2{,}52\degg$ & $5{,}30\degg$ \\
\textbf{EKF + LDR terkalibrasi} & $\bm{0{,}18\degg}$ & $0{,}39\degg$ \\
EKF + LDR $3\%$ tidak terkalibrasi & $\bm{5{,}63\degg}$ & $6{,}10\degg$ \\
EKF + LDR terkalibrasi + arus & $0{,}18\degg$ & $0{,}39\degg$ \\
EKF + arus saja & $2{,}24\degg$ & $5{,}17\degg$ \\
\bottomrule
\end{longtable}

Tapis pulih dari lepas langkah $7{,}2\degg$ dalam sekitar $350$~detik \textbf{tanpa pernah diberi tahu}
bahwa lepas langkah terjadi.

\textbf{Tiga kesimpulan.} (i)~LDR terkalibrasi memperbaiki empat belas kali. (ii)~\textbf{LDR tidak
terkalibrasi lebih buruk daripada tidak memakai LDR sama sekali}: bias LDR tidak dapat dibedakan dari salah
pasang mekanik (tidak teramati secara terpisah), sehingga tapis menyerapnya lalu mengarahkan panel keliru.
(iii)~Arus hampir tidak menambah apa pun, sesuai $\partial I/\partial \bm{b} \approx 0$ di dekat puncak.

\subsection{Keyhole sebagai kovarians}

\begin{longtable}{l r r}
\toprule
\textbf{Keadaan} & $\partial e_y/\partial b_{\mathrm{az}}$ & $\sigma_{\mathrm{az}}$ \textbf{dilaporkan} \\
\midrule
Matahari elevasi $45\degg$ & $0{,}00091$ & $0{,}67\degg$ \\
Matahari elevasi $85\degg$ & $0{,}00019$ & $\bm{2{,}65\degg}$ \\
\bottomrule
\end{longtable}

\subsection{Dua pengaman yang wajib}

Keduanya ditemukan ketika loop tertutup menyimpang di Gazebo.

\begin{enumerate}[leftmargin=1.4em]
\item \textbf{LDR hanya dipakai saat matahari terjangkau.} Bila zenith melewati $45\degg$, kemiringan
terpotong dan LDR membaca galat besar yang bukan berasal dari simpangan sendi.
\item \textbf{Anti-penyimpangan pada gerbang pencilan.} Bila taksiran telanjur salah, inovasi membesar,
gerbang menolak semua pengukuran, dan tapis terkunci selamanya (terukur: NIS tetap $82$ sampai $90$,
taksiran macet di $-6{,}47\degg$). Setelah sepuluh penolakan beruntun, $\bm{P}$ digembungkan.
\end{enumerate}

\subsection{Pelajaran tentang matriks kovarians pengukuran}

Matriks $\bm{R}$ untuk pengukuran arus \textbf{wajib memuat galat model}, bukan hanya derau sensor:
\begin{equation}
R_I = \sigma_{\mathrm{ADC}}^2 + \left( \kappa_{\mathrm{model}} \hat{I} \right)^2,
\qquad \kappa_{\mathrm{model}} = 0{,}03 .
\end{equation}
Mengisinya dengan derau ADC $0{,}2$~mA saja membuat $\bm{K} = \bm{P}\bm{H}^T\bm{S}^{-1}$ ribuan kali terlalu
besar, dan \textbf{galat pointing melonjak dari $0{,}18\degg$ menjadi $20{,}5\degg$}.

\subsection{Kejujuran tapis}

Tapis masih \emph{terlalu percaya diri} sekitar $3{,}4$ kali pada azimuth: $\sigma$ dilaporkan $0{,}23\degg$
sedangkan sebaran galat sebenarnya $0{,}79\degg$. Sebabnya $\bm{Q}$ berbentuk jalan acak halus sedangkan
lepas langkah datang sebagai lompatan.

%=====================================================================
\section{Arsitektur Cyber-Physical System}

\begin{Verbatim}[fontsize=\small,frame=single]
 +--------------------------------------------------------------------+
 |                        LAPISAN CYBER                               |
 |  sun_command_node  -->  profil trapesium + ZVD  <--  ekf_node      |
 |  (pvlib SPA, IK)        (100 Hz)                     (taksir b, P) |
 |  sensor_node  --> /tracker/diag (5 Hz), /ldr/state        ^        |
 |  wind_node    --> angin nyata berarah (Open-Meteo)        |        |
 +-----------------------------+------------------------------+------+
                               | perintah posisi (rad)        |
 +-----------------------------v------------------------------+------+
 | LAPISAN ANTARMUKA : ros_gz_bridge                                  |
 |   kuantisasi microstep 0,1125 deg   <-- Q(.)                       |
 |   kejenuhan torsi 3 N m             <-- Sat(.)                     |
 +-----------------------------+------------------------------+------+
                               | tau(t)         /joint_states (1 kHz) |
 +-----------------------------v------------------------------+------+
 | LAPISAN FISIK : Gazebo DART, langkah 1 ms                          |
 |   M(q) qdd + C(q,qd) qd + g(q) + F(qd) + tau_w = tau               |
 |   gangguan: angin berarah, defleksi gravitasi, lepas langkah       |
 +--------------------------------------------------------------------+
\end{Verbatim}

\begin{longtable}{p{5.2cm} p{5.6cm} p{4.0cm}}
\toprule
\textbf{Kerangka rujukan} & \textbf{Di sistem ini} & \textbf{Status} \\
\midrule
\endhead
Plant dinamika kontinu & Gazebo DART, langkah 1 ms & berjalan \\
Kuantisasi $Q(\cdot)$ & pembulatan microstep $0{,}1125\degg$ & berjalan \\
Kejenuhan $\mathrm{Sat}(\cdot)$ & batas torsi $3$ N\,m & berjalan \\
Algoritma orbit astronomis & pvlib SPA & tervalidasi $0{,}011\degg$ \\
Pengendali diskret & node ROS 2, 100 Hz & berjalan \\
Penaksir keadaan (EKF) & \texttt{ekf\_node} & berjalan \\
Enkoder kuadratur & \textbf{ada, di rotor}, tertutup di dalam driver & tidak terbaca host; sisi beban tetap ditaksir EKF \\
Latensi jaringan, jitter & \textbf{belum dimodelkan} & kekurangan \\
Serangan FDI, NTP poisoning & \textbf{belum dimodelkan} & kekurangan \\
Event-triggered control & \textbf{belum ada} & kekurangan \\
\bottomrule
\end{longtable}

\textbf{Catatan tentang jam.} Seluruh node memakai \texttt{use\_sim\_time}. Tanpa itu, profil gerak dan
input shaper berjalan pada jam dinding sedangkan plant pada jam simulasi (Gazebo sekitar $85\%$
real-time), sehingga lewatan menjadi tidak dapat diulang: terukur $7{,}6\%$ sampai $0{,}35\%$ pada uji
yang sama.

\subsection{Model keadaan hibrid}

\model{Sistem sampled-data dengan tundaan dan kuantisasi}

Kerangka rujukan menuliskan bentuk hibrid yang relevan dan belum diterapkan di sini:
\begin{align}
\tilde{\bm{y}}[k] &= Q\big( \bm{x}(kT_s - \tau_{d,s}) \big) + \bm{v}[k], \\
\bm{u}_c[k] &= \mathcal{C}\big( \tilde{\bm{y}}[k], \bm{q}_d[k], \hat{\bm{\theta}}[k] \big), \\
\bm{\tau}(t) &= \mathrm{Sat}\big( \bm{u}_c[k - d_a] + \bm{a}_k \big),
\quad t \in [kT_s + \tau_d,\ (k+1)T_s + \tau_d),
\end{align}
dengan $d_a$ kehilangan paket, $\tau_d$ tundaan total, dan $\bm{a}_k$ penyuntikan data palsu. Kuantisasi
$Q$ dan kejenuhan $\mathrm{Sat}$ sudah ada; $\tau_{d,s}$, $d_a$, dan $\bm{a}_k$ belum.

%=====================================================================
\section{Lapisan Digital Twin}

\model{Taksonomi Kritzinger}

\begin{longtable}{l l l}
\toprule
\textbf{Tingkat} & \textbf{Aliran data} & \textbf{Posisi} \\
\midrule
\textbf{Digital Model} & tidak ada, manual & \textbf{di sini} \\
Digital Shadow & otomatis, fisik $\rightarrow$ digital & belum \\
Digital Twin & otomatis, dua arah & belum \\
\bottomrule
\end{longtable}

Yang ada adalah Digital Model \emph{tervalidasi silang tiga jalur}: nilai eigen analitik, model Python dua
sumbu, dan Gazebo saling cocok. Yang belum ada adalah hubungan ke alat fisik.

\textbf{Jembatannya pendek.} Node \texttt{sensor\_node} menghasilkan data sensor dan menerbitkannya ke
\texttt{/tracker/diag}; \texttt{ekf\_node} memakannya dari topik itu. Mengganti \texttt{sensor\_node} dengan
node pembaca perangkat keras yang menerbitkan pesan sama membuat seluruh rantai di belakangnya berjalan
tanpa perubahan.

Usulan kerangka rujukan yang belum ada: identifikasi parameter $\bm{\theta}$ secara daring, optimasi
multi-objektif prediktif
$\max \sum_k [w_1 P_{\mathrm{gen}} - w_2 P_{\mathrm{act}} - w_3 D_{\mathrm{fatigue}}]$,
indeks kesehatan dari hanyutan gesekan, dan deteksi \emph{backlash}. Dua usulan lain tidak berlaku:
\emph{intelligent backtracking} ditujukan untuk ladang surya banyak panel (alat ini satu panel), dan
\emph{predictive aero-elastic wind stow} tidak mendesak karena margin angin sebelas kali.

\textbf{Catatan untuk pekerjaan lanjut.} Biaya kelelahan struktur
$D_{\mathrm{fatigue}} = \lVert \bm{\tau}(k) - \bm{\tau}(k-1) \rVert^2$ memberi pembenaran kuantitatif bagi
peredam ZVD, yang selama ini hanya dapat dibenarkan secara kualitatif karena rugi energinya $0{,}007\%$.

%=====================================================================
\section{Ringkasan Hasil Terukur}

\begin{longtable}{p{6.4cm} p{4.6cm} p{4.0cm}}
\toprule
\textbf{Besaran} & \textbf{Nilai} & \textbf{Cara diperoleh} \\
\midrule
\endhead
Posisi matahari & selisih $< 0{,}011\degg$ & SPA vs astropy, SolTrack, Pysolar \\
Kinematika normal panel & selisih $0{,}14$--$0{,}34\degg$ & Gazebo TF vs Pers.~\eqref{eq:fk} \\
Defleksi gravitasi $q_2 = 0$ & $0{,}161\degg$ vs $0{,}155\degg$ & Gazebo vs Pers.~\eqref{eq:grav} \\
Frekuensi alami azimuth & $4{,}79$ prediksi, $5{,}19$ terukur Hz & nilai eigen vs Gazebo 1 kHz \\
Lewatan langkah $1\degg$ & $93{,}3\%$ teori, $90\%$ terukur & analitik vs Gazebo \\
ZVD menekan lewatan terburuk & $19{,}35\degg \rightarrow 2{,}30\degg$ & Gazebo, lima ulangan \\
Rugi energi dari osilasi & $0{,}007\%$ & model dua sumbu 1 kHz \\
Energi harian 7 Okt & $431{,}0$ Wh vs $362{,}3$ Wh datar & model dioda tunggal \\
Margin angin terburuk 11 tahun & $9{,}2\%$ torsi tahan & Open-Meteo + Pers.~\eqref{eq:wind} \\
EKF dengan LDR terkalibrasi & $2{,}52\degg \rightarrow 0{,}18\degg$ & simulasi hari penuh \\
Pemulihan lepas langkah $7{,}2\degg$ & $350$ detik & simulasi hari penuh \\
\bottomrule
\end{longtable}

%=====================================================================
\section{Kekurangan dan Parameter yang Masih Diasumsikan}
\label{sec:asumsi}

\textbf{Belum dimodelkan:} suku Coriolis (besarnya sudah dihitung, $3{,}4\%$ saat slew), latensi dan jitter
jaringan, serangan siber, \emph{event-triggered control}, optimasi multi-objektif, \emph{backlash} kopling,
dan mode getar rangka.

\textbf{Asumsi yang paling menentukan:}

\begin{longtable}{p{5.4cm} p{3.0cm} p{6.4cm}}
\toprule
\textbf{Parameter} & \textbf{Nilai} & \textbf{Cara mengukur} \\
\midrule
\endhead
Redaman sendi $b$ & $0{,}2$ N\,m\,s/rad & \textbf{uji peluruhan bebas} \\
Gesekan Coulomb $\tau_c$ & $0{,}05$ N\,m & torsi minimum penggerak \\
Kekakuan poros $k$ & $317$ dan $1600$ N\,m/rad & uji peluruhan bebas \\
Eksponen LDR $\gamma$ & $0{,}7$ & kurva lux-tegangan di meja \\
Setengah jenuh $L_0$ & $20000$ lux & kurva yang sama \\
Ketidakcocokan LDR & $3\%$ & prosedur A saat mendung \\
Batas kemiringan & $45\degg$ & gerakkan sampai menumbuk \\
Arah hadap rel $\psi$ & $90\degg$ & kompas \\
Massa motor & $1{,}2$ kg & timbang atau lembar data \\
Koefisien angin $C_n$ & $1{,}2$ & tabel beban angin atau CFD \\
Pusat tekanan $e_c$, $e_s$ & $0{,}0668$, $0{,}02$ m & aerodinamika pelat datar \\
\bottomrule
\end{longtable}

\textbf{Uji peluruhan bebas} adalah percobaan tunggal paling bernilai: dari satu rekaman langkah dapat
diperoleh $\omega_n$ (dari frekuensi ayunan), $\zeta$ lalu $b$ (dari rasio puncak berurutan melalui
$\zeta = -\ln(\text{rasio})/2\pi$), dan $\tau_c$ (dari simpangan sisa yang tidak kembali ke nol).

\textbf{Dua klaim yang ditarik dalam pengerjaan:} batas angin tarik-lepas $9{,}6$~m/s (lengan momen tebakan;
rentang sebenarnya $6{,}5$ sampai $39$~m/s) dan pergeseran puncak arus $+2{,}6\degg$ (irradiance pilihan
tangan tidak konsisten dengan $G_h = I_b\sin e + I_d$).

%=====================================================================
\begin{thebibliography}{99}
\small
\bibitem{spa} I. Reda dan A. Andreas, \emph{Solar Position Algorithm for Solar Radiation Applications},
NREL/TP-560-34302. Dipakai lewat \texttt{pvlib}.
\bibitem{ineichen} P. Ineichen dan R. Perez, \emph{A new airmass independent formulation for the Linke
turbidity coefficient}, Solar Energy, 2002.
\bibitem{kasten} F. Kasten dan G. Czeplak, \emph{Solar and terrestrial radiation dependent on the amount and
type of cloud}, Solar Energy, 1980.
\bibitem{erbs} D. G. Erbs, S. A. Klein, dan J. A. Duffie, \emph{Estimation of the diffuse radiation fraction},
Solar Energy, 1982.
\bibitem{liu} B. Y. H. Liu dan R. C. Jordan, model transposisi isotropik untuk irradiance bidang miring.
\bibitem{ashrae} ASHRAE, model \emph{incidence angle modifier}; lihat juga Souka dan Safwat.
\bibitem{desoto} W. De Soto, S. A. Klein, dan W. A. Beckman, \emph{Improvement and validation of a model for
photovoltaic array performance}, Solar Energy, 2006.
\bibitem{singer} N. C. Singer dan W. P. Seering, \emph{Preshaping command inputs to reduce system vibration},
ASME J. Dynamic Systems, Measurement and Control, 1990.
\bibitem{spong} M. W. Spong, S. Hutchinson, dan M. Vidyasagar, \emph{Robot Modeling and Control}
(persamaan manipulator dan simbol Christoffel).
\bibitem{barshalom} Y. Bar-Shalom, X. R. Li, dan T. Kirubarajan, \emph{Estimation with Applications to
Tracking and Navigation} (EKF, bentuk Joseph, uji NIS).
\bibitem{kritzinger} W. Kritzinger dkk., \emph{Digital Twin in manufacturing: A categorical literature
review and classification}, IFAC-PapersOnLine, 2018.
\bibitem{openmeteo} Open-Meteo Historical Weather API, reanalisis ERA5 per jam.
\end{thebibliography}

\end{document}
